\documentclass[11pt]{article}
\usepackage[margin=1in]{geometry}
\usepackage[T1]{fontenc}
\usepackage[utf8]{inputenc}
\usepackage{lmodern}
\usepackage[varqu,scaled=0.95]{zi4}
\usepackage{amsmath,amssymb,amsthm,mathtools}
\usepackage[dvipsnames]{xcolor}
\usepackage{fvextra}
\usepackage{booktabs,enumitem,microtype,needspace}
\usepackage[colorlinks,linkcolor=MidnightBlue,citecolor=MidnightBlue,urlcolor=MidnightBlue]{hyperref}
\usepackage[capitalise,nameinlink]{cleveref}
\newtheorem{theorem}{Theorem}[section]
\newtheorem{lemma}[theorem]{Lemma}
\theoremstyle{definition}
\newtheorem{definition}[theorem]{Definition}
\theoremstyle{remark}
\newtheorem{remark}[theorem]{Remark}
\newcommand{\leanname}[1]{\texttt{\detokenize{#1}}}
\newenvironment{leanblock}[1]
  {\par\smallskip\Needspace{4\baselineskip}\noindent{\footnotesize\textsf{#1}}\par\nopagebreak\vspace{2pt}}
  {\par\medskip}
\fvset{breaklines=true,breakanywhere=true,fontsize=\footnotesize,frame=leftline,
  framerule=0.8pt,rulecolor=\color{gray!60},xleftmargin=4pt,framesep=6pt}
\usepackage{longtable}
\newcommand{\unverified}{\quad\textup{\small[not machine-verified]}}

\newcommand{\R}{\mathbb{R}}
\newcommand{\N}{\mathbb{N}}
\newcommand{\Pn}{\mathcal{P}_n}
\newcommand{\one}{\mathbf{1}}
\newcommand{\wn}[1]{\lVert #1\rVert_{w,1}}
\newcommand{\hn}[1]{\lVert #1\rVert_{w}}
\newcommand{\ip}[2]{\langle #1,#2\rangle_{w}}
\newcommand{\lnorm}[1]{\lVert #1\rVert_{1}}

\title{Metric Markov cotype two of $\ell_1$ with constant $3024$: an exposition of a machine-checked proof}
\date{}
\begin{document}
\maketitle
\begin{abstract}
We prove that $\ell_1$ has metric Markov cotype two with $C^2=3024$, that is, $N_2(\ell_1)\le 12\sqrt{21}$. For every stationary pair $(\pi,A)$ on $[n]$, every $t\ge 1$ and every $x_1,\dots,x_n\in\ell_1$, there are points $y_i\in\ell_1$ whose combined fidelity and smoothness cost is at most $3024$ times the Ces\`aro-averaged energy of the $x_i$. The argument has three parts. First, the points are encoded by weighted cut coordinates. Second, the cubic map $r\mapsto 3r^2-2r^3$ is applied coordinatewise to resolvent averages of the binary encodings, and its energy is bounded by a quartic potential that telescopes against the resolvent $p\sum_s q^sA^s$ with $p=1/(t+1)$. Third, the geometric (resolvent) average of endpoint costs is compared with the Ces\`aro average, in the spirit of Mendel--Naor. The main theorem and all lemmas below are machine-verified. The probabilistic martingale of the source argument is replaced in the formal proof by an equivalent algebraic telescoping.
\end{abstract}

\begin{center}\small
\begin{tabular}{@{}ll@{}}
\toprule
Source & openai/math@fd4aeeb, family 332, OpenAIMath\_332\_MetricMarkovCotypeL1 \\
Status & proof machine-verified (Lean 4, Mathlib) \\
Lemmas & 49 \\
Text & written by claude-opus-5-5, checked against the Lean source \\
Run & \texttt{srv\_20261009\_010221\_OpenAIMath\_332\_MetricMarkovCotypeL1} \\
\bottomrule
\end{tabular}
\end{center}

\tableofcontents

\section{Setting}

\paragraph{Index sets and matrices.}
Throughout, $\N=\{0,1,2,\dots\}$. For $n\in\N$ let $[n]=\{1,\dots,n\}$ with its natural order; $[n]$ is empty when $n=0$. Vectors are indexed by $[n]$ and matrices by $[n]\times[n]$. For a matrix $A$ we write $a_{ij}=A_{ij}$, we write $\delta_{ij}$ for the Kronecker delta, and we set $A^0=I$ and $A^{s+1}=A^sA$. We write $\Pn$ for the set of all $2^n$ subsets of $[n]$, including $\emptyset$ and $[n]$. For $i\in[n]$ the binary vector $\chi_i\in\{0,1\}^{\Pn}$ is defined by $\chi_i(B)=1$ if $i\in B$ and $\chi_i(B)=0$ otherwise.

For $s,t\in\N$ with $t\ge1$, $\lfloor s/t\rfloor$ and $s\bmod t$ are the quotient and remainder of Euclidean division, so that $s=t\lfloor s/t\rfloor+(s\bmod t)$. For $t=0$ we use the conventions $\lfloor s/0\rfloor=0$ and $s\bmod 0=s$, and also $1/0=0$. These conventions only affect degenerate cases, which are pointed out where they occur.

\paragraph{Series.}
For a real sequence $(c_s)_{s\in\N}$ we write $\sum_{s\ge0}c_s$ for the sum of the series when the series is summable (equivalently, absolutely convergent), and $0$ otherwise. Whenever a lemma asserts convergence, it says so explicitly. In every place where this convention is used, the proof shows that the series converges.

\begin{definition}[The space $\ell_1$]
$\ell_1$ is the real vector space of sequences $f=(f(k))_{k\in\N}$ with $\sum_k|f(k)|^p<\infty$ for $p=1$. It carries pointwise operations and the norm $\lnorm{f}=\bigl(\sum_k|f(k)|^p\bigr)^{1/p}$ with $p=1$. This makes $\ell_1$ a real Banach space. Indexing by $\N$ instead of $\{1,2,\dots\}$ is immaterial, since the index shift is a linear isometry onto.
\end{definition}

\begin{definition}[Probability vectors, stochastic and reversible matrices]
Let $n\in\N$.
\begin{itemize}
\item A vector $\pi\in\R^n$ is a \emph{probability vector} if $\pi_i\ge0$ for all $i$ and $\sum_i\pi_i=1$. Zero entries are allowed.
\item A matrix $A\in\R^{n\times n}$ is \emph{stochastic} if $a_{ij}\ge0$ for all $i,j$ and $\sum_ja_{ij}=1$ for all $i$.
\item $A$ is \emph{reversible} with respect to $\pi\in\R^n$ if $\pi_ia_{ij}=\pi_ja_{ji}$ for all $i,j$.
\item $\pi$ is \emph{stationary} for $A$, written $\pi A=\pi$, if $\sum_i\pi_ia_{ij}=\pi_j$ for all $j$.
\end{itemize}
\end{definition}

\paragraph{Resolvent parameters.}
For $t\in\N$ put
\[
p=p_t=\frac1{t+1},\qquad q=q_t=\frac t{t+1}.
\]
Then $p>0$, $0\le q<1$, $p+q=1$, $(t+1)p=1$ and $(t+1)q=t$. For a stochastic matrix $A$, the \emph{resolvent matrix} $G$ is the matrix with entries
\[
G_{ij}=\sum_{s\ge0}p\,q^s(A^s)_{ij}.
\]
These series converge by \cref{lem:mc_resolvent_summable}.

\paragraph{Weighted seminorms and the cubic map.}
Let $\iota$ be a finite set and $w\in\R^\iota$ with $w_b\ge0$ for all $b$. For $u,v\in\R^\iota$ put
\[
\wn{u}=\sum_{b}w_b|u_b|,\qquad \ip uv=\sum_bw_bu_bv_b,\qquad \hn u^2=\ip uu,\qquad \hn u^4=(\hn u^2)^2 .
\]
Let $\phi(r)=3r^2-2r^3$ for $r\in\R$. For $u\in\R^\iota$ let $\Phi(u)\in\R^\iota$ be defined by $\Phi(u)_b=\phi(u_b)$. We write $Q=[0,1]^\iota$ for the cube and $\{0,1\}^\iota$ for its set of vertices.

\begin{theorem}[Metric Markov cotype two of $\ell_1$]\label{thm:main}
Let $n,t\in\N$ with $t\ge1$. Let $\pi$ be a probability vector on $[n]$, let $A$ be a stochastic matrix on $[n]$ that is reversible with respect to $\pi$, and let $x_1,\dots,x_n\in\ell_1$. Then there exist $y_1,\dots,y_n\in\ell_1$ such that
\[
\sum_i\pi_i\lnorm{x_i-y_i}^2+t\sum_{i,j}\pi_ia_{ij}\lnorm{y_i-y_j}^2
\le 3024\sum_{i,j}\pi_i\Bigl(\frac1t\sum_{s=1}^t(A^s)_{ij}\Bigr)\lnorm{x_i-x_j}^2 .
\]
\end{theorem}

This is exactly the statement of the source, with $C^2=3024$, so that $N_2(\ell_1)\le\sqrt{3024}=12\sqrt{21}$. The only formal difference is that sequences in $\ell_1$ are indexed from $0$, which is an isometric relabelling. As in the source, reversibility enters only through stationarity.

\section{Lemmas}

\subsection{Stationarity and powers of stochastic matrices}

\begin{lemma}[Reversibility implies stationarity]\label{lem:mc_stationary}
Let $\pi\in\R^n$ and let $A$ be stochastic and reversible with respect to $\pi$. Then $\sum_i\pi_ia_{ij}=\pi_j$ for every $j$.
\end{lemma}
\begin{proof}
By detailed balance and the row-sum condition,
\[
\sum_i\pi_ia_{ij}=\sum_i\pi_ja_{ji}=\pi_j\sum_ia_{ji}=\pi_j. \qedhere
\]
\end{proof}

\begin{lemma}[Powers of stochastic matrices]\label{lem:mc_pow_stochastic}
If $A$ is stochastic, then $A^s$ is stochastic for every $s\in\N$.
\end{lemma}
\begin{proof}
We argue by induction on $s$. The identity matrix is stochastic. If $A^s$ is stochastic, then $(A^{s+1})_{ij}=\sum_k(A^s)_{ik}a_{kj}\ge0$. Moreover
\[
\sum_j\sum_k(A^s)_{ik}a_{kj}=\sum_k(A^s)_{ik}\sum_ja_{kj}=\sum_k(A^s)_{ik}=1. \qedhere
\]
\end{proof}

\begin{lemma}[Stationarity of powers]\label{lem:mc_pow_stationary}
Let $\pi\in\R^n$ and $A\in\R^{n\times n}$ satisfy $\pi A=\pi$. Then $\sum_i\pi_i(A^s)_{ij}=\pi_j$ for all $s\in\N$ and all $j$.
\end{lemma}
\begin{proof}
We argue by induction on $s$; the case $s=0$ is trivial. Using $A^{s+1}=A^sA$, exchanging finite sums and applying the induction hypothesis,
\[
\sum_i\pi_i\sum_k(A^s)_{ik}a_{kj}=\sum_k\Bigl(\sum_i\pi_i(A^s)_{ik}\Bigr)a_{kj}=\sum_k\pi_ka_{kj}=\pi_j. \qedhere
\]
\end{proof}

\subsection{The resolvent}

\begin{lemma}[Resolvent series converge]\label{lem:mc_resolvent_summable}
Let $t\in\N$, let $A$ be stochastic and let $i,j\in[n]$. Then the series $\sum_{s\ge0}p\,q^s(A^s)_{ij}$ is summable.
\end{lemma}
\begin{proof}
First, $0\le(A^s)_{ij}\le1$ for all $s$. This follows by induction from $(A^{s+1})_{ij}=\sum_ka_{ik}(A^s)_{kj}$: the sum is nonnegative, and it is at most $\sum_ka_{ik}=1$. Hence $0\le p\,q^s(A^s)_{ij}\le p\,q^s$. Since $0\le q<1$, the geometric series $\sum_s p\,q^s$ converges, and the claim follows by comparison.
\end{proof}

\begin{lemma}[Basic properties of the resolvent]\label{lem:mc_resolvent_basic}
Let $t\in\N$ and $\pi\in\R^n$. Let $A$ be stochastic with $\pi A=\pi$, and let $G$ be the resolvent matrix, $G_{ij}=\sum_{s\ge0}p\,q^s(A^s)_{ij}$. Then:
\begin{enumerate}
\item $G_{ij}\ge0$ for all $i,j$;
\item $\sum_jG_{ij}=1$ for all $i$;
\item $\sum_i\pi_iG_{ij}=\pi_j$ for all $j$.
\end{enumerate}
\end{lemma}
\begin{proof}
Since $1-q=p$, we have $\sum_{s\ge0}p\,q^s=p/(1-q)=1$. All series $\sum_sp\,q^s(A^s)_{ij}$ converge by \cref{lem:mc_resolvent_summable}, so finite sums and constant multiples may be taken termwise.
\begin{enumerate}
\item Each term is nonnegative by \cref{lem:mc_pow_stochastic}.
\item By \cref{lem:mc_pow_stochastic}, $\sum_jG_{ij}=\sum_sp\,q^s\sum_j(A^s)_{ij}=\sum_sp\,q^s=1$.
\item By \cref{lem:mc_pow_stationary}, $\sum_i\pi_iG_{ij}=\sum_sp\,q^s\sum_i\pi_i(A^s)_{ij}=\sum_sp\,q^s\pi_j=\pi_j$. \qedhere
\end{enumerate}
\end{proof}

\begin{lemma}[First-step identities]\label{lem:mc_resolvent_first_step}
Let $t\in\N$, let $A$ be stochastic and let $G$ be its resolvent matrix. Then for all $i,j$,
\[
G_{ij}=p\,\delta_{ij}+q\sum_ka_{ik}G_{kj}
\qquad\text{and}\qquad
G_{ij}=p\,\delta_{ij}+q\sum_kG_{ik}a_{kj}.
\]
\end{lemma}
\begin{proof}
All series involved converge by \cref{lem:mc_resolvent_summable}. Hence
\[
\sum_ka_{ik}G_{kj}=\sum_{s\ge0}p\,q^s\sum_ka_{ik}(A^s)_{kj}=\sum_{s\ge0}p\,q^s(A^{s+1})_{ij}.
\]
Splitting off the term $s=0$ of $G_{ij}$, which equals $p\,\delta_{ij}$, gives
\[
G_{ij}=p\,\delta_{ij}+\sum_{s\ge0}p\,q^{s+1}(A^{s+1})_{ij}=p\,\delta_{ij}+q\sum_ka_{ik}G_{kj}.
\]
The second identity follows in the same way from $A^{s+1}=A^sA$.
\end{proof}

\begin{lemma}[Resolvent averages of binary vectors]\label{lem:mc_harmonic_point}
Let $t\in\N$ and $A,G\in\R^{n\times n}$. Assume $G_{ij}\ge0$, $\sum_jG_{ij}=1$ and $G_{ij}=p\,\delta_{ij}+q\sum_ka_{ik}G_{kj}$ for all $i,j$. Let $\iota$ be any set and let $z_1,\dots,z_n:\iota\to\{0,1\}$. Put $h_i(b)=\sum_jG_{ij}z_j(b)$. Then for all $i\in[n]$ and $b\in\iota$:
\begin{enumerate}
\item $0\le h_i(b)\le1$;
\item $h_i(b)=p\,z_i(b)+q\sum_ka_{ik}h_k(b)$.
\end{enumerate}
\end{lemma}
\begin{proof}
Part~(1) holds because $h_i(b)$ is a convex combination of numbers in $\{0,1\}$. For part~(2), multiply the hypothesis $G_{ij}=p\,\delta_{ij}+q\sum_ka_{ik}G_{kj}$ by $z_j(b)$, sum over $j$, and exchange the finite sums over $j$ and $k$:
\[
h_i(b)=p\,z_i(b)+q\sum_ka_{ik}\sum_jG_{kj}z_j(b). \qedhere
\]
\end{proof}

\subsection{The cubic map and the quartic potential}

\begin{lemma}[Scalar increment bound]\label{lem:mc_phi_scalar}
For $r,r'\in[0,1]$ and $d\in\{0,1\}$,
\[
|\phi(r')-\phi(r)|\le6|r-d|\,|r'-r|+3(r'-r)^2 .
\]
\end{lemma}
\begin{proof}
Let $\delta=r'-r$. Factoring the difference of squares and of cubes gives
\[
\phi(r')-\phi(r)=\delta\,K,\qquad K=3(r'+r)-2\bigl(r'^2+r'r+r^2\bigr).
\]
It therefore suffices to show
\[
|K|\le6|r-d|+3|\delta|,
\]
since multiplying by $|\delta|$ then gives the claim. In the formal proof this bound is obtained by case analysis on $d\in\{0,1\}$ and on the sign of $\delta$; in each case it is a polynomial inequality for $(r,r')\in[0,1]^2$, discharged by nonlinear arithmetic. A direct way to see it is as follows. Substituting $r'=r+\delta$ gives $K=6r(1-r)+\delta(3-4r-2r')$.
\begin{itemize}
\item Since $4r+2r'\in[0,6]$, we have $|3-4r-2r'|\le3$.
\item We have $0\le6r(1-r)\le6|r-d|$. If $d=0$, then $|r-d|=r$ and $1-r\le1$. If $d=1$, then $|r-d|=1-r$ and $r\le1$.
\end{itemize}
The triangle inequality gives the bound on $|K|$.
\end{proof}

\begin{lemma}[Weighted increment bound]\label{lem:mc_flat_increment}
Let $\iota$ be finite, $w\in\R^\iota$ with $w\ge0$, $u,u'\in[0,1]^\iota$ and $e\in\{0,1\}^\iota$. Then
\[
\wn{\Phi(u')-\Phi(u)}^2\le72\,\hn{u-e}^2\,\hn{u'-u}^2+18\,\hn{u'-u}^4 .
\]
\end{lemma}
\begin{proof}
Put
\[
X=\wn{\Phi(u')-\Phi(u)},\qquad S=\sum_bw_b|u_b-e_b|\,|u'_b-u_b|,\qquad P=\hn{u'-u}^2.
\]

\emph{Linear bound.} The formal proof does not cite \cref{lem:mc_phi_scalar}. It re-establishes the same scalar bound inline: for $r,r'\in[0,1]$ and $d\in\{0,1\}$,
\[
|\phi(r')-\phi(r)|\le6|r-d|\,|r'-r|+3(r'-r)^2 .
\]
This inline argument uses the expansion $\phi(r')-\phi(r)=6r(1-r)\,\delta+\delta^2(3-6r-2\delta)$ with $\delta=r'-r$, together with the estimates $|3-6r-2\delta|=|3-4r-2r'|\le3$ and $0\le6r(1-r)\le6|r-d|$. Applying this bound in each coordinate with $d=e_b$, multiplying by $w_b\ge0$ and summing gives $0\le X\le6S+3P$.

\emph{Cauchy--Schwarz.} Writing $w_b=(\sqrt{w_b})^2$, the Cauchy--Schwarz inequality gives $S^2\le\hn{u-e}^2P$.

\emph{Squaring.} From $(6S-3P)^2\ge0$ we get $36SP\le36S^2+9P^2$. Therefore
\[
X^2\le(6S+3P)^2\le72S^2+18P^2\le72\,\hn{u-e}^2P+18P^2. \qedhere
\]
\end{proof}

\begin{lemma}[Flatness against the quartic potential]\label{lem:mc_flat}
Under the hypotheses of \cref{lem:mc_flat_increment},
\[
\wn{\Phi(u')-\Phi(u)}^2\le108\Bigl(\hn{u'-e}^4-\hn{u-e}^4-4\,\hn{u-e}^2\,\ip{u-e}{u'-u}\Bigr).
\]
\end{lemma}
\begin{proof}
Put $S=\hn{u-e}^2$, $P=\hn{u'-u}^2$ and $I=\ip{u-e}{u'-u}$. Then $\hn{u'-e}^2=S+2I+P$, so the bracket equals
\[
R=(S+2I+P)^2-S^2-4SI=2SP+(2I+P)^2 .
\]
By Cauchy--Schwarz (again with $\sqrt{w_b}$), $I^2\le SP$. Hence
\[
P^2\le2(2I+P)^2+2(2I)^2\le2(2I+P)^2+8SP .
\]
Combining this with \cref{lem:mc_flat_increment},
\[
\wn{\Phi(u')-\Phi(u)}^2\le72SP+18P^2\le216SP+36(2I+P)^2\le108R. \qedhere
\]
\end{proof}

\begin{lemma}[One averaged step]\label{lem:mc_one_step}
Let $\iota$ be finite, $w\in\R^\iota$ with $w\ge0$ and $e\in\{0,1\}^\iota$. Let $p,q\ge0$ with $p+q=1$, and let $a\in\R^n$ with $a_k\ge0$ and $\sum_ka_k=1$. Let $u,v,h_1,\dots,h_n\in[0,1]^\iota$ satisfy
\[
u_b=p\,v_b+q\sum_ka_kh_k(b)\quad\text{for all } b .
\]
Then
\begin{multline*}
p\,\wn{\Phi(v)-\Phi(u)}^2+q\sum_ka_k\wn{\Phi(h_k)-\Phi(u)}^2\\
\le108\Bigl(p\,\hn{v-e}^4+q\sum_ka_k\hn{h_k-e}^4-\hn{u-e}^4\Bigr).
\end{multline*}
\end{lemma}
\begin{proof}
Apply \cref{lem:mc_flat} to the pair $(u,v)$ and to each pair $(u,h_k)$. Multiply these inequalities by $p$ and by $q\,a_k$ respectively, which are nonnegative, and add them.

\emph{Constant terms.} Since $p+q\sum_ka_k=1$, the terms $-\hn{u-e}^4$ combine to $-\hn{u-e}^4$.

\emph{Linear terms.} These carry the common factor $-4\hn{u-e}^2$ and the cross term
\[
p\,\ip{u-e}{v-u}+q\sum_ka_k\ip{u-e}{h_k-u}
=\sum_bw_b(u_b-e_b)\Bigl(p\,v_b+q\sum_ka_kh_k(b)-(p+q)u_b\Bigr).
\]
This vanishes by the hypothesis on $u$ and $p+q=1$.
\end{proof}

\begin{lemma}[Aggregation against the resolvent]\label{lem:mc_aggregate}
Let $t\in\N$ and $\pi\in\R^n$ with $\pi_i\ge0$. Let $A,G\in\R^{n\times n}$ satisfy:
\begin{itemize}
\item $G_{ij}\ge0$ and $G_{ij}=p\,\delta_{ij}+q\sum_kG_{ik}a_{kj}$ for all $i,j$;
\item $\sum_i\pi_iG_{ij}=\pi_j$ for all $j$.
\end{itemize}
Let $\beta\in\R^n$ and $E,F,H\in\R^{n\times n}$ with $H_{ij}\ge0$ for all $i,j$. Assume that for all $i,j\in[n]$,
\[
p\,\beta_j+q\sum_ka_{jk}E_{jk}\le108\Bigl(p\,F_{ij}+q\sum_ka_{jk}H_{ik}-H_{ij}\Bigr).
\]
Then
\[
\sum_j\pi_j\beta_j+t\sum_j\sum_k\pi_ja_{jk}E_{jk}\le108\sum_i\sum_j\pi_iG_{ij}F_{ij}.
\]
\end{lemma}
\begin{proof}
Multiply the hypothesis at $(i,j)$ by $\pi_iG_{ij}\ge0$ and sum over $i$ and $j$.

\emph{Left side.} Using $\sum_i\pi_iG_{ij}=\pi_j$, it becomes
\[
\sum_j\pi_j\Bigl(p\beta_j+q\sum_ka_{jk}E_{jk}\Bigr).
\]

\emph{Right side.} Fix $i$. By the hypothesis on $G$, $q\sum_jG_{ij}a_{jk}=G_{ik}-p\,\delta_{ik}$. Therefore
\[
\sum_jG_{ij}\,q\sum_ka_{jk}H_{ik}=\sum_k\bigl(G_{ik}-p\,\delta_{ik}\bigr)H_{ik}=\sum_kG_{ik}H_{ik}-p\,H_{ii}.
\]
The $i$-th contribution to the right side is thus
\[
108\,\pi_i\Bigl(p\sum_jG_{ij}F_{ij}-p\,H_{ii}\Bigr)\le108\,p\,\pi_i\sum_jG_{ij}F_{ij},
\]
since $\pi_ip\,H_{ii}\ge0$.

\emph{Conclusion.} Summing over $i$ gives
\[
p\sum_j\pi_j\beta_j+q\sum_{j,k}\pi_ja_{jk}E_{jk}\le108\,p\sum_{i,j}\pi_iG_{ij}F_{ij}.
\]
Multiplying by $t+1>0$ and using $(t+1)p=1$ and $(t+1)q=t$ gives the claim.
\end{proof}

\begin{lemma}[Vertices are fixed]\label{lem:mc_phi_vertex}
If $d\in\{0,1\}$, then $\phi(d)=d$.
\end{lemma}
\begin{proof}
$\phi(0)=0$ and $\phi(1)=3-2=1$.
\end{proof}

\begin{lemma}[One step for the resolvent averages]\label{lem:mc_cube_step}
Let $t\in\N$ and let $A$ be stochastic. Let $\iota$ be finite, $w\in\R^\iota$ with $w\ge0$, and $z_1,\dots,z_n\in\{0,1\}^\iota$. Let $G\in\R^{n\times n}$ satisfy $G_{ij}\ge0$, $\sum_jG_{ij}=1$ and $G_{ij}=p\,\delta_{ij}+q\sum_ka_{ik}G_{kj}$ for all $i,j$. Put $h_k=\sum_lG_{kl}z_l\in\R^\iota$. Then for all $i_0,j\in[n]$,
\begin{multline*}
p\,\wn{\Phi(z_j)-\Phi(h_j)}^2+q\sum_ka_{jk}\wn{\Phi(h_k)-\Phi(h_j)}^2\\
\le108\Bigl(p\,\hn{z_j-z_{i_0}}^4+q\sum_ka_{jk}\hn{h_k-z_{i_0}}^4-\hn{h_j-z_{i_0}}^4\Bigr).
\end{multline*}
\end{lemma}
\begin{proof}
We first derive two facts directly from the hypotheses on $G$. These are the conclusions of \cref{lem:mc_harmonic_point}, but the formal proof re-derives them rather than citing it.
\begin{itemize}
\item \emph{$h_k\in[0,1]^\iota$ for all $k$.} For each $b$, $h_k(b)=\sum_lG_{kl}z_l(b)$ with $G_{kl}\ge0$, $\sum_lG_{kl}=1$ and $z_l(b)\in\{0,1\}$. Hence $0\le h_k(b)\le\sum_lG_{kl}=1$.
\item \emph{The first-step identity $h_j=p\,z_j+q\sum_ka_{jk}h_k$.} Multiply $G_{jl}=p\,\delta_{jl}+q\sum_ka_{jk}G_{kl}$ by $z_l(b)$, sum over $l$, and exchange the finite sums over $l$ and $k$.
\end{itemize}
Moreover $z_j\in\{0,1\}^\iota\subseteq[0,1]^\iota$, $p,q\ge0$ and $p+q=1$, and the row $(a_{jk})_k$ is nonnegative with sum $1$. Apply \cref{lem:mc_one_step} with $e=z_{i_0}$, $a_k=a_{jk}$, $u=h_j$, $v=z_j$ and the vectors $h_k$.
\end{proof}

\begin{lemma}[Energy bound, raw form]\label{lem:mc_cube_energy_raw}
Let $t\in\N$, let $\pi\in\R^n$ with $\pi_i\ge0$ and let $A$ be stochastic. Let $\iota$ be finite, $w\in\R^\iota$ with $w\ge0$, and $z_1,\dots,z_n\in\{0,1\}^\iota$. Let $G\in\R^{n\times n}$ satisfy:
\begin{itemize}
\item $G_{ij}\ge0$ and $\sum_jG_{ij}=1$;
\item both identities $G_{ij}=p\,\delta_{ij}+q\sum_ka_{ik}G_{kj}$ and $G_{ij}=p\,\delta_{ij}+q\sum_kG_{ik}a_{kj}$;
\item $\sum_i\pi_iG_{ij}=\pi_j$.
\end{itemize}
Let $h_k=\sum_lG_{kl}z_l$. Then
\[
\sum_j\pi_j\wn{\Phi(z_j)-\Phi(h_j)}^2+t\sum_j\sum_k\pi_ja_{jk}\wn{\Phi(h_k)-\Phi(h_j)}^2
\le108\sum_i\sum_j\pi_iG_{ij}\hn{z_j-z_i}^4 .
\]
\end{lemma}
\begin{proof}
Apply \cref{lem:mc_aggregate} with
\[
\beta_j=\wn{\Phi(z_j)-\Phi(h_j)}^2,\quad E_{jk}=\wn{\Phi(h_k)-\Phi(h_j)}^2,\quad F_{ij}=\hn{z_j-z_i}^4,\quad H_{ij}=\hn{h_j-z_i}^4\ge0 .
\]
Its hypothesis at $(i,j)$ is exactly \cref{lem:mc_cube_step} with $i_0=i$.
\end{proof}

\begin{lemma}[Energy bound]\label{lem:mc_cube_energy}
Assume the hypotheses of \cref{lem:mc_cube_energy_raw}, except that $\pi$ is now a probability vector. Then
\[
\sum_i\pi_i\wn{z_i-\Phi(h_i)}^2+t\sum_i\sum_j\pi_ia_{ij}\wn{\Phi(h_i)-\Phi(h_j)}^2
\le108\sum_i\sum_j\pi_iG_{ij}\hn{z_i-z_j}^4 .
\]
\end{lemma}
\begin{proof}
This is \cref{lem:mc_cube_energy_raw}, which only uses $\pi\ge0$, rewritten in two ways. First, $\Phi(z_i)=z_i$ by \cref{lem:mc_phi_vertex}. Second, the seminorms are invariant under $v\mapsto-v$.
\end{proof}

\subsection{Cut decomposition in $\ell_1$}

\begin{lemma}[Cuts of a monotone vector]\label{lem:mc_cut_real_mono}
Let $n\in\N$ and let $a\in\R^n$ be nondecreasing, that is, $i\le j$ implies $a_i\le a_j$. Then there exist subsets $c_1,\dots,c_n\subseteq[n]$ and reals $g_1,\dots,g_n,v_0$ such that:
\begin{enumerate}
\item $g_r\ge0$ for all $r$;
\item $\sum_rg_r\le2\sum_i|a_i|$;
\item $|v_0|\le\sum_i|a_i|$;
\item $a_i=v_0+\sum_r\one[i\in c_r]\,g_r$ for all $i$;
\item $|a_i-a_j|=\sum_r\bigl|\one[i\in c_r]-\one[j\in c_r]\bigr|\,g_r$ for all $i,j$.
\end{enumerate}
\end{lemma}
\begin{proof}
If $n=0$, everything is empty and all claims are trivial. Otherwise define
\[
c_r=\{i: i>r\},\qquad g_r=a_{r+1}-a_r\ (r<n),\qquad g_n=0,\qquad v_0=a_1 .
\]
\begin{itemize}
\item Monotonicity gives $g_r\ge0$.
\item By telescoping, $\sum_{r<m}g_r=a_m-a_1$ for every $m\in[n]$.
\item Since $\one[i\in c_r]\,g_r=\one[r<i]\,g_r$, the reconstruction~(4) follows.
\item We have $\sum_rg_r=a_n-a_1\le|a_n|+|a_1|\le2\sum_i|a_i|$, and $|v_0|=|a_1|\le\sum_i|a_i|$.
\end{itemize}
For~(5), let $i\le j$. A three-case check ($r<i$, then $i\le r<j$, then $r\ge j$) gives
\[
\bigl|\one[i\in c_r]-\one[j\in c_r]\bigr|g_r=\one[r<j]g_r-\one[r<i]g_r .
\]
Summing over $r$ gives $(a_j-a_1)-(a_i-a_1)=|a_i-a_j|$. The case $i>j$ follows by symmetry.
\end{proof}

\begin{lemma}[Cuts of an arbitrary vector]\label{lem:mc_cut_real_pointwise}
The conclusion of \cref{lem:mc_cut_real_mono} holds for every $a\in\R^n$, without the monotonicity assumption.
\end{lemma}
\begin{proof}
Choose a permutation $\sigma$ of $[n]$ such that $a\circ\sigma$ is nondecreasing; such a $\sigma$ exists by sorting. Apply \cref{lem:mc_cut_real_mono} to $a\circ\sigma$, obtaining $c'_r,g_r,v_0$, and set $c_r=\sigma(c'_r)$. Then $\one[i\in c_r]=\one[\sigma^{-1}(i)\in c'_r]$. Hence~(4) and~(5) for $a$ follow from those for $a\circ\sigma$ evaluated at $\sigma^{-1}(i)$ and $\sigma^{-1}(j)$. Finally, $\sum_i|a_{\sigma(i)}|=\sum_i|a_i|$.
\end{proof}

\begin{lemma}[Uniform cut decomposition of $\R^n$]\label{lem:mc_cut_real}
Let $n\in\N$. There exist functions $g:\Pn\times\R^n\to\R$, $(B,a)\mapsto g_B(a)$, and $v:\R^n\to\R$ such that for every $a\in\R^n$:
\begin{enumerate}
\item $0\le g_B(a)\le2\sum_i|a_i|$ for all $B\in\Pn$;
\item $|v(a)|\le\sum_i|a_i|$;
\item $a_i=v(a)+\sum_{B\in\Pn}\one[i\in B]\,g_B(a)$ for all $i$;
\item $|a_i-a_j|=\sum_{B\in\Pn}\bigl|\one[i\in B]-\one[j\in B]\bigr|\,g_B(a)$ for all $i,j$.
\end{enumerate}
\end{lemma}
\begin{proof}
For each $a$ choose data $c_r(a),g_r(a),v_0(a)$ as in \cref{lem:mc_cut_real_pointwise}. Put $v(a)=v_0(a)$ and
\[
g_B(a)=\sum_{r:\,c_r(a)=B}g_r(a).
\]
For any function $f$ on $\Pn$ we have $\sum_Bf(B)\,g_B(a)=\sum_rf(c_r(a))\,g_r(a)$. Applying this to $f(B)=\one[i\in B]$ and to $f(B)=|\one[i\in B]-\one[j\in B]|$ transfers~(3) and~(4) from \cref{lem:mc_cut_real_pointwise}. Finally, $0\le g_B(a)\le\sum_rg_r(a)\le2\sum_i|a_i|$.
\end{proof}

\begin{lemma}[The $\ell_1$ norm]\label{lem:mc_L1_norm}
For $f\in\ell_1$, the sequence $(|f(k)|)_k$ is summable and $\lnorm f=\sum_k|f(k)|$.
\end{lemma}
\begin{proof}
This is the definition of membership in $\ell_p$ and of the $\ell_p$ norm, specialised to $p=1$.
\end{proof}

\begin{lemma}[Domination gives membership in $\ell_1$]\label{lem:mc_L1_of_bound}
Let $f,g:\N\to\R$ with $g$ summable and $|f(k)|\le g(k)$ for all $k$. Then there exists $F\in\ell_1$ with $F(k)=f(k)$ for all $k$.
\end{lemma}
\begin{proof}
By comparison, $\sum_k|f(k)|<\infty$, so $f$ itself belongs to $\ell_1$.
\end{proof}

\begin{lemma}[Coordinates of affine combinations]\label{lem:mc_L1_affine_apply}
Let $\iota$ be finite, $V\in\ell_1$, $b_B\in\ell_1$ for $B\in\iota$, and $u\in\R^\iota$. Then for every $k\in\N$,
\[
\Bigl(V+\sum_Bu_Bb_B\Bigr)(k)=V(k)+\sum_Bu_Bb_B(k).
\]
\end{lemma}
\begin{proof}
The vector operations on $\ell_1$ are pointwise.
\end{proof}

\begin{lemma}[Affine maps are weighted contractions]\label{lem:mc_L1_affine_contr}
In the setting of \cref{lem:mc_L1_affine_apply}, for all $u,u'\in\R^\iota$,
\[
\Bigl\|\Bigl(V+\sum_Bu_Bb_B\Bigr)-\Bigl(V+\sum_Bu'_Bb_B\Bigr)\Bigr\|_1\le\sum_B\lnorm{b_B}\,|u_B-u'_B| .
\]
\end{lemma}
\begin{proof}
The difference is $\sum_B(u_B-u'_B)b_B$. The claim follows from the triangle inequality and homogeneity of the norm.
\end{proof}

\begin{lemma}[Distances from coordinatewise cut identities]\label{lem:mc_L1_dist_eq}
Let $\iota$ be finite, let $x_1,\dots,x_n\in\ell_1$, and let $b_B\in\ell_1$ for $B\in\iota$ with $b_B(k)\ge0$ for all $B,k$. Let $z_1,\dots,z_n\in\R^\iota$, and assume that for all $i,j,k$,
\[
|x_i(k)-x_j(k)|=\sum_B|z_i(B)-z_j(B)|\,b_B(k).
\]
Then $\lnorm{x_i-x_j}=\sum_B\lnorm{b_B}\,|z_i(B)-z_j(B)|$ for all $i,j$.
\end{lemma}
\begin{proof}
By \cref{lem:mc_L1_norm}, each $b_B$ is summable, and $\lnorm{b_B}=\sum_kb_B(k)$ because $b_B\ge0$. Using \cref{lem:mc_L1_norm} for $x_i-x_j$, the hypothesis, and the fact that a finite sum of convergent series may be summed termwise,
\[
\lnorm{x_i-x_j}=\sum_k\sum_B|z_i(B)-z_j(B)|\,b_B(k)=\sum_B|z_i(B)-z_j(B)|\sum_kb_B(k). \qedhere
\]
\end{proof}

\begin{lemma}[Cut coordinates in $\ell_1$]\label{lem:mc_cut_coords}
Let $x_1,\dots,x_n\in\ell_1$. There exist $b_B\in\ell_1$ for $B\in\Pn$ and $V\in\ell_1$ such that, for all $B$, $i$, $j$, $k$:
\begin{enumerate}
\item $b_B(k)\ge0$;
\item $x_i(k)=V(k)+\sum_B\one[i\in B]\,b_B(k)$;
\item $|x_i(k)-x_j(k)|=\sum_B\bigl|\one[i\in B]-\one[j\in B]\bigr|\,b_B(k)$.
\end{enumerate}
\end{lemma}
\begin{proof}
Take $g$ and $v$ from \cref{lem:mc_cut_real}. For each $k$ let $a^{(k)}=(x_1(k),\dots,x_n(k))$, and define the sequences $\beta_B(k)=g_B(a^{(k)})$ and $\nu(k)=v(a^{(k)})$.

By \cref{lem:mc_L1_norm}, the sequence $k\mapsto\sum_i|x_i(k)|$ is a finite sum of summable sequences, hence summable. By \cref{lem:mc_cut_real},
\[
|\beta_B(k)|\le2\sum_i|x_i(k)|\qquad\text{and}\qquad|\nu(k)|\le\sum_i|x_i(k)| .
\]
So \cref{lem:mc_L1_of_bound} yields $b_B=\beta_B\in\ell_1$ and $V=\nu\in\ell_1$. Properties~(1)--(3) are \cref{lem:mc_cut_real} applied to $a^{(k)}$.
\end{proof}

\begin{lemma}[Cut representation with a contractive reconstruction]\label{lem:mc_cut_L1}
Let $x_1,\dots,x_n\in\ell_1$. There exist $w\in\R^{\Pn}$ and a map $T:\R^{\Pn}\to\ell_1$ such that:
\begin{enumerate}
\item $w_B\ge0$ for all $B$;
\item $T(\chi_i)=x_i$ for all $i$;
\item $\lnorm{x_i-x_j}=\sum_Bw_B|\chi_i(B)-\chi_j(B)|$ for all $i,j$;
\item $\lnorm{T(u)-T(u')}\le\sum_Bw_B|u_B-u'_B|$ for all $u,u'\in\R^{\Pn}$.
\end{enumerate}
\end{lemma}
\begin{proof}
Take $b_B$ and $V$ from \cref{lem:mc_cut_coords}, and put $w_B=\lnorm{b_B}$ and $T(u)=V+\sum_Bu_Bb_B$; thus $T$ is affine.
\begin{itemize}
\item By \cref{lem:mc_L1_affine_apply} and \cref{lem:mc_cut_coords}(2), $T(\chi_i)(k)=x_i(k)$ for every $k$, which gives~(2).
\item Property~(3) is \cref{lem:mc_L1_dist_eq} with $z_i=\chi_i$, using \cref{lem:mc_cut_coords}(1) and~(3).
\item Property~(4) is \cref{lem:mc_L1_affine_contr}. \qedhere
\end{itemize}
\end{proof}

\begin{lemma}[Transfer from the cube to $\ell_1$]\label{lem:mc_transfer}
Let $t\in\N$, let $\pi\in\R^n$ with $\pi\ge0$, and let $A\in\R^{n\times n}$ with $a_{ij}\ge0$. Let $x_1,\dots,x_n\in\ell_1$, let $\iota$ be finite, $w\in\R^\iota$ with $w\ge0$, and $z_1,\dots,z_n\in\{0,1\}^\iota$. Let $T:\R^\iota\to\ell_1$ satisfy:
\begin{itemize}
\item $T(z_i)=x_i$;
\item $\lnorm{x_i-x_j}=\wn{z_i-z_j}$;
\item $\lnorm{T(u)-T(u')}\le\wn{u-u'}$ for all $u,u'$.
\end{itemize}
Let $G\in\R^{n\times n}$ be arbitrary and let $c_1,\dots,c_n\in\R^\iota$ satisfy
\[
\sum_i\pi_i\wn{z_i-c_i}^2+t\sum_{i,j}\pi_ia_{ij}\wn{c_i-c_j}^2\le108\sum_{i,j}\pi_iG_{ij}\hn{z_i-z_j}^4 .
\]
Then there exist $y_1,\dots,y_n\in\ell_1$ with
\[
\sum_i\pi_i\lnorm{x_i-y_i}^2+t\sum_{i,j}\pi_ia_{ij}\lnorm{y_i-y_j}^2\le108\sum_{i,j}\pi_iG_{ij}\lnorm{x_i-x_j}^2 .
\]
\end{lemma}
\begin{proof}
Put $y_i=T(c_i)$. Then
\[
\lnorm{x_i-y_i}=\lnorm{T(z_i)-T(c_i)}\le\wn{z_i-c_i}\qquad\text{and}\qquad\lnorm{y_i-y_j}\le\wn{c_i-c_j}.
\]
Both sides are nonnegative, so these inequalities survive squaring. Multiplying by $\pi_i\ge0$ and by $t\,\pi_ia_{ij}\ge0$ and summing shows that the left side is at most the left side of the hypothesis.

For binary numbers $|\alpha-\beta|=(\alpha-\beta)^2$. Hence $\hn{z_i-z_j}^2=\wn{z_i-z_j}=\lnorm{x_i-x_j}$, so $\hn{z_i-z_j}^4=\lnorm{x_i-x_j}^2$, and the right sides coincide.
\end{proof}

\subsection{Geometric versus Ces\`aro averages}

\begin{lemma}[Weighted Minkowski inequality]\label{lem:mc_minkowski}
Let $\alpha$ be finite and $\mu,a,b\in\R^\alpha$ with $\mu\ge0$. Then
\[
\sqrt{\sum_k\mu_k(a_k+b_k)^2}\le\sqrt{\sum_k\mu_ka_k^2}+\sqrt{\sum_k\mu_kb_k^2}.
\]
\end{lemma}
\begin{proof}
This is the triangle inequality for the Euclidean norm, applied to the vectors $(\sqrt{\mu_k}a_k)_k$ and $(\sqrt{\mu_k}b_k)_k$.
\end{proof}

\begin{lemma}[Weighted triangle inequality]\label{lem:mc_weighted_tri}
Let $\alpha$ be finite and $\mu,a,b,c\in\R^\alpha$ with $\mu\ge0$, $c\ge0$ and $c_k\le a_k+b_k$ for all $k$. Then
\[
\sqrt{\sum_k\mu_kc_k^2}\le\sqrt{\sum_k\mu_ka_k^2}+\sqrt{\sum_k\mu_kb_k^2}.
\]
\end{lemma}
\begin{proof}
Let $\alpha_0$ and $\beta_0$ denote the two square roots on the right. Since $0\le c_k\le|a_k|+|b_k|$, we have $c_k^2\le a_k^2+2|a_k||b_k|+b_k^2$. By Cauchy--Schwarz with weights $\sqrt{\mu_k}$, $\sum_k\mu_k|a_k||b_k|\le\alpha_0\beta_0$. Hence
\[
\sum_k\mu_kc_k^2\le\alpha_0^2+2\alpha_0\beta_0+\beta_0^2=(\alpha_0+\beta_0)^2 .
\]
Take square roots.
\end{proof}

\begin{lemma}[Expanding a product of two matrices]\label{lem:mc_triple_mul}
For $\pi\in\R^n$ and $P,Q,d\in\R^{n\times n}$,
\[
\sum_{i,j}\pi_i(PQ)_{ij}d_{ij}=\sum_{(i,k,j)\in[n]^3}\pi_iP_{ik}Q_{kj}d_{ij}.
\]
\end{lemma}
\begin{proof}
Expand $(PQ)_{ij}=\sum_kP_{ik}Q_{kj}$ and rearrange the finite sum.
\end{proof}

\begin{lemma}[Summing out the last index]\label{lem:mc_triple_first}
Let $\pi\in\R^n$ and $P,Q,d\in\R^{n\times n}$ with $\sum_jQ_{kj}=1$ for all $k$. Then
\[
\sum_{(i,k,j)}\pi_iP_{ik}Q_{kj}d_{ik}=\sum_{i,k}\pi_iP_{ik}d_{ik}.
\]
\end{lemma}
\begin{proof}
Sum over $j$ first and use $\sum_jQ_{kj}=1$.
\end{proof}

\begin{lemma}[Summing out the first index]\label{lem:mc_triple_second}
Let $\pi\in\R^n$ and $P,Q,d\in\R^{n\times n}$ with $\sum_i\pi_iP_{ik}=\pi_k$ for all $k$. Then
\[
\sum_{(i,k,j)}\pi_iP_{ik}Q_{kj}d_{kj}=\sum_{k,j}\pi_kQ_{kj}d_{kj}.
\]
\end{lemma}
\begin{proof}
Sum over $i$ first and use $\sum_i\pi_iP_{ik}=\pi_k$.
\end{proof}

\begin{lemma}[Square-root subadditivity of the endpoint energy]\label{lem:mc_D_subadd}
Let $\pi$ be a probability vector, $A$ stochastic with $\pi A=\pi$, and $x_1,\dots,x_n\in\ell_1$. Put
\[
D(s)=\sum_{i,j}\pi_i(A^s)_{ij}\lnorm{x_i-x_j}^2 .
\]
Then $\sqrt{D(r+s)}\le\sqrt{D(r)}+\sqrt{D(s)}$ for all $r,s\in\N$.
\end{lemma}
\begin{proof}
Let $\mu(i,k,j)=\pi_i(A^r)_{ik}(A^s)_{kj}$; it is nonnegative by \cref{lem:mc_pow_stochastic}. We express all three quantities as $\mu$-weighted sums.
\begin{itemize}
\item Since $A^{r+s}=A^rA^s$, \cref{lem:mc_triple_mul} gives $D(r+s)=\sum\mu(i,k,j)\lnorm{x_i-x_j}^2$.
\item Since the rows of $A^s$ sum to $1$ (\cref{lem:mc_pow_stochastic}), \cref{lem:mc_triple_first} gives $D(r)=\sum\mu(i,k,j)\lnorm{x_i-x_k}^2$.
\item Since $\pi A^r=\pi$ (\cref{lem:mc_pow_stationary}), \cref{lem:mc_triple_second} gives $D(s)=\sum\mu(i,k,j)\lnorm{x_k-x_j}^2$.
\end{itemize}
The triangle inequality $\lnorm{x_i-x_j}\le\lnorm{x_i-x_k}+\lnorm{x_k-x_j}$ and \cref{lem:mc_weighted_tri} conclude.
\end{proof}

\begin{lemma}[Bound on $D(t)$]\label{lem:mc_cesaro_Dt}
Let $t\ge1$ and let $D:\N\to\R$ satisfy $D(0)=0$, $D\ge0$ and $\sqrt{D(r+s)}\le\sqrt{D(r)}+\sqrt{D(s)}$ for all $r,s$. Then
\[
D(t)\le4\cdot\frac1t\sum_{s=1}^tD(s).
\]
\end{lemma}
\begin{proof}
For $1\le r\le t$, squaring $\sqrt{D(t)}\le\sqrt{D(r)}+\sqrt{D(t-r)}$ and using $(\alpha+\beta)^2\le2\alpha^2+2\beta^2$ gives $D(t)\le2D(r)+2D(t-r)$. Sum over $r=1,\dots,t$. As $r$ runs over $1,\dots,t$, the index $t-r$ runs over $0,\dots,t-1$, and $D(0)=0$. Hence
\[
t\,D(t)\le2\sum_{r=1}^tD(r)+2\sum_{s=1}^{t-1}D(s)\le4\sum_{s=1}^tD(s),
\]
using $D\ge0$ in the last step. Divide by $t$.
\end{proof}

\begin{lemma}[Block decomposition]\label{lem:mc_cesaro_blocks}
Let $t\in\N$ and let $D:\N\to\R$ satisfy $D\ge0$ and $\sqrt{D(r+s)}\le\sqrt{D(r)}+\sqrt{D(s)}$ for all $r,s$. Then for all $j,r\in\N$,
\[
D(jt+r)\le2j^2D(t)+2D(r).
\]
\end{lemma}
\begin{proof}
By induction on $j$, using $(j+1)t+r=(jt+r)+t$ and subadditivity, we get $\sqrt{D(jt+r)}\le j\sqrt{D(t)}+\sqrt{D(r)}$. Squaring (both sides are nonnegative) and using $(\alpha+\beta)^2\le2\alpha^2+2\beta^2$ gives the claim.
\end{proof}

\begin{lemma}[Second moment of the geometric series]\label{lem:mc_hasSum_sq_geom}
For $0\le r<1$, the series $\sum_{s\ge0}s^2r^s$ converges, and its sum is $r(1+r)/(1-r)^3$.
\end{lemma}
\begin{proof}
We have $s^2=2\binom{s+2}{2}-3s-2$. We also have the standard sums
\[
\sum_s\binom{s+2}2r^s=\frac1{(1-r)^3},\qquad\sum_ssr^s=\frac r{(1-r)^2},\qquad\sum_sr^s=\frac1{1-r},
\]
all of which converge. Combining them linearly,
\[
\frac{2}{(1-r)^3}-\frac{3r}{(1-r)^2}-\frac{2}{1-r}=\frac{2-3r+3r^2-2+4r-2r^2}{(1-r)^3}=\frac{r(1+r)}{(1-r)^3}. \qedhere
\]
\end{proof}

\begin{lemma}[Floor bound]\label{lem:mc_floor_div_sq_le}
For $t\ge1$ and $s\in\N$, $\lfloor s/t\rfloor^2\le s^2/t^2$.
\end{lemma}
\begin{proof}
We have $t\lfloor s/t\rfloor\le s$, so $0\le\lfloor s/t\rfloor\le s/t$. Square.
\end{proof}

\begin{lemma}[Geometric second moment of the block index]\label{lem:mc_geom_sq_div}
Let $t\ge1$. Then the series $\sum_{s\ge0}p\,q^s\lfloor s/t\rfloor^2$ converges and its sum is at most $3$.
\end{lemma}
\begin{proof}
By \cref{lem:mc_floor_div_sq_le}, the terms are nonnegative and bounded by $\frac{p}{t^2}s^2q^s$. By \cref{lem:mc_hasSum_sq_geom}, these bounds are summable with sum
\[
\frac p{t^2}\cdot\frac{q(1+q)}{(1-q)^3}.
\]
Using $1-q=1/(t+1)$ and $q=t/(t+1)$, this equals $(2t+1)/t\le3$.
\end{proof}

\begin{lemma}[Geometric weight of a residue class]\label{lem:mc_geom_mod_coeff}
Let $t\ge1$ and $r,N\in\N$. Then
\[
\sum_{\substack{0\le s<N\\ s\bmod t=r}}p\,q^s\le\frac2t .
\]
\end{lemma}
\begin{proof}
\emph{The factor $q^t$.} We have $q(1+1/t)=1$, and by Bernoulli's inequality $(1+1/t)^t\ge1+t\cdot\frac1t=2$. Hence $q^t\le1/2$.

\emph{Termwise bound.} For $s$ in the sum, $q^s=q^{s\bmod t}(q^t)^{\lfloor s/t\rfloor}\le(1/2)^{\lfloor s/t\rfloor}$, because $0\le q\le1$.

\emph{Reindexing.} The map $s\mapsto\lfloor s/t\rfloor$ is injective on the residue class, since $s=(s\bmod t)+t\lfloor s/t\rfloor$. It maps $\{0,\dots,N-1\}$ into itself, since $\lfloor s/t\rfloor\le s$. Therefore the sum is at most
\[
p\sum_{k<N}2^{-k}=p\,(2-2\cdot2^{-N})\le2p\le\frac2t. \qedhere
\]
\end{proof}

\begin{lemma}[Geometric average of a periodic function]\label{lem:mc_geom_mod}
Let $t\ge1$ and $f:\N\to\R$ with $f\ge0$. Then the series $\sum_{s\ge0}p\,q^sf(s\bmod t)$ converges, and its sum is at most $\frac2t\sum_{r=0}^{t-1}f(r)$.
\end{lemma}
\begin{proof}
\emph{Convergence.} The terms are nonnegative and bounded by $p\,q^s\sum_{r<t}f(r)$, which is summable. So the series converges.

\emph{Bound.} For every $N$, group the partial sum by the residue $r=s\bmod t\in\{0,\dots,t-1\}$ and apply \cref{lem:mc_geom_mod_coeff}:
\[
\sum_{s<N}p\,q^sf(s\bmod t)=\sum_{r<t}f(r)\sum_{\substack{s<N\\ s\bmod t=r}}p\,q^s\le\frac2t\sum_{r<t}f(r).
\]
A series of nonnegative terms whose partial sums are bounded by a constant has sum at most that constant.
\end{proof}

\begin{lemma}[Geometric averages of bounded sequences]\label{lem:mc_geom_summable_bdd}
Let $t\in\N$ and let $f:\N\to\R$ satisfy $f\ge0$ and $f\le C$ for some constant $C$. Then the series $\sum_{s\ge0}p\,q^sf(s)$ converges.
\end{lemma}
\begin{proof}
We have $0\le p\,q^sf(s)\le C\,p\,q^s$, and $\sum_sq^s$ converges because $0\le q<1$.
\end{proof}

\begin{lemma}[Euclidean-division bound]\label{lem:mc_D_divmod}
Let $t\in\N$ and let $D:\N\to\R$ satisfy $D\ge0$ and $\sqrt{D(r+s)}\le\sqrt{D(r)}+\sqrt{D(s)}$ for all $r,s$. Then for every $s\in\N$,
\[
D(s)\le2\lfloor s/t\rfloor^2D(t)+2D(s\bmod t).
\]
For $t=0$ this is read with the conventions $\lfloor s/0\rfloor=0$ and $s\bmod0=s$.
\end{lemma}
\begin{proof}
If $t=0$, the claim reads $D(s)\le2D(s)$, which holds since $D\ge0$.

Let $t\ge1$, and write $m=\lfloor s/t\rfloor$ and $\rho=s\bmod t$, so that $s=mt+\rho$. The formal proof does not cite \cref{lem:mc_cesaro_blocks}; it reproduces the same bound by the following route.
\begin{enumerate}
\item \emph{Two-term bound.} For all $a,b\in\N$, squaring $\sqrt{D(a+b)}\le\sqrt{D(a)}+\sqrt{D(b)}$ (both sides are nonnegative) and using $(\alpha+\beta)^2\le2\alpha^2+2\beta^2$ gives $D(a+b)\le2D(a)+2D(b)$.
\item \emph{Multiples of $t$.} For $m\ge1$, we have $\sqrt{D(mt)}\le m\sqrt{D(t)}$, by induction on $m$. The case $m=1$ is trivial. For the induction step, $\sqrt{D((m+1)t)}\le\sqrt{D(mt)}+\sqrt{D(t)}\le(m+1)\sqrt{D(t)}$. Squaring gives $D(mt)\le m^2D(t)$.
\item \emph{The case $m=0$.} Then $s=\rho$, and $D(s)=D(\rho)\le2D(\rho)=2\cdot0^2\cdot D(t)+2D(\rho)$, since $D\ge0$.
\item \emph{The case $m\ge1$.} By steps~1 and~2,
\[
D(s)=D(mt+\rho)\le2D(mt)+2D(\rho)\le2m^2D(t)+2D(\rho).
\]
\end{enumerate}
This is the bound of \cref{lem:mc_cesaro_blocks} with $j=m$ and $r=\rho$. \qedhere
\end{proof}

\begin{lemma}[Splitting the geometric average]\label{lem:mc_geom_cesaro_pre}
Let $t\ge1$ and let $D:\N\to\R$ be nonnegative, bounded, and square-root subadditive as in \cref{lem:mc_D_divmod}. Then
\[
\sum_{s\ge0}p\,q^sD(s)\le2D(t)\sum_{s\ge0}p\,q^s\lfloor s/t\rfloor^2+2\sum_{s\ge0}p\,q^sD(s\bmod t).
\]
\end{lemma}
\begin{proof}
\emph{Convergence.} The series on the left and the last series converge, since $D$ is bounded and nonnegative (compare with $C\,p\,q^s$). The middle series converges because $0\le\lfloor s/t\rfloor^2q^s\le s^2q^s$, which is summable by \cref{lem:mc_hasSum_sq_geom}.

\emph{Termwise bound.} Multiplying \cref{lem:mc_D_divmod} by $p\,q^s\ge0$ bounds the series termwise. Comparison of convergent series and linearity of the sum give the claim.
\end{proof}

\begin{lemma}[Shifting a Ces\`aro window]\label{lem:mc_range_le_Icc}
Let $t\in\N$ and let $D:\N\to\R$ satisfy $D(0)=0$ and $D\ge0$. Then
\[
\sum_{r=0}^{t-1}D(r)\le\sum_{s=1}^tD(s).
\]
\end{lemma}
\begin{proof}
For $t=0$ both sides vanish. For $t\ge1$, the left side equals $\sum_{r=1}^{t-1}D(r)$ because $D(0)=0$, and this is at most the same sum plus $D(t)\ge0$.
\end{proof}

\begin{lemma}[Geometric versus Ces\`aro averages]\label{lem:mc_geom_cesaro}
Let $t\ge1$ and let $D:\N\to\R$ satisfy $D(0)=0$, $D\ge0$, $D$ bounded, and $\sqrt{D(r+s)}\le\sqrt{D(r)}+\sqrt{D(s)}$ for all $r,s$. Then
\[
\sum_{s\ge0}p\,q^sD(s)\le28\cdot\frac1t\sum_{s=1}^tD(s).
\]
\end{lemma}
\begin{proof}
Let $W=\frac1t\sum_{s=1}^tD(s)$. By \cref{lem:mc_geom_cesaro_pre}, the left side is at most $2D(t)\cdot S_2+2M$, where $S_2=\sum_sp\,q^s\lfloor s/t\rfloor^2$ and $M=\sum_sp\,q^sD(s\bmod t)$. We bound each piece.
\begin{itemize}
\item $S_2\le3$ by \cref{lem:mc_geom_sq_div}; since $D(t)\ge0$, this gives $2D(t)S_2\le6D(t)$.
\item $M\le\frac2t\sum_{r<t}D(r)\le2W$ by \cref{lem:mc_geom_mod} and \cref{lem:mc_range_le_Icc}.
\item $D(t)\le4W$ by \cref{lem:mc_cesaro_Dt}.
\end{itemize}
Altogether the left side is at most $6\cdot4W+2\cdot2W=28W$.
\end{proof}

\begin{lemma}[Resolvent energy as a geometric average]\label{lem:mc_resolvent_D_eq}
Let $t\in\N$, $\pi\in\R^n$, $A$ stochastic, $x_1,\dots,x_n\in\ell_1$, and let $G$ be the resolvent matrix of $A$. With $D(s)=\sum_{i,j}\pi_i(A^s)_{ij}\lnorm{x_i-x_j}^2$,
\[
\sum_{i,j}\pi_iG_{ij}\lnorm{x_i-x_j}^2=\sum_{s\ge0}p\,q^sD(s).
\]
\end{lemma}
\begin{proof}
Each series $G_{ij}=\sum_sp\,q^s(A^s)_{ij}$ converges by \cref{lem:mc_resolvent_summable}. Multiply it by $\pi_i\lnorm{x_i-x_j}^2$ and add these finitely many convergent series. This shows that $\sum_sp\,q^sD(s)$ converges, and that its sum is the left side.
\end{proof}

\begin{lemma}[Properties of the endpoint energy]\label{lem:mc_D_props}
Let $\pi$ be a probability vector, $A$ stochastic, $x_1,\dots,x_n\in\ell_1$, and $D$ as in \cref{lem:mc_resolvent_D_eq}. Then:
\begin{enumerate}
\item $D(0)=0$;
\item $D(s)\ge0$ for all $s$;
\item there is a constant $C$ with $D(s)\le C$ for all $s$.
\end{enumerate}
\end{lemma}
\begin{proof}
\begin{enumerate}
\item Since $A^0=I$, $D(0)=\sum_i\pi_i\lnorm{x_i-x_i}^2=0$.
\item Nonnegativity follows from $\pi\ge0$ and $(A^s)_{ij}\ge0$ (\cref{lem:mc_pow_stochastic}).
\item Each entry of $A^s$ is nonnegative and at most its row sum $1$, and $\pi_i\le\sum_k\pi_k=1$. Hence $D(s)\le C:=\sum_{i,j}\lnorm{x_i-x_j}^2$. \qedhere
\end{enumerate}
\end{proof}

\begin{lemma}[Ces\`aro energy]\label{lem:mc_cesaro_rhs_eq}
Let $t\in\N$, $\pi\in\R^n$, $A\in\R^{n\times n}$, $x_1,\dots,x_n\in\ell_1$, and $D$ as above. Then
\[
\sum_{i,j}\pi_i\Bigl(\frac1t\sum_{s=1}^t(A^s)_{ij}\Bigr)\lnorm{x_i-x_j}^2=\frac1t\sum_{s=1}^tD(s).
\]
For $t=0$ both sides vanish, since the inner sums are empty.
\end{lemma}
\begin{proof}
Pull out the factor $1/t$ and exchange the finite sums over $s$ and over $(i,j)$.
\end{proof}

\begin{lemma}[Comparison of resolvent and Ces\`aro energies]\label{lem:mc_comparison}
Let $t\ge1$, let $\pi$ be a probability vector, let $A$ be stochastic with $\pi A=\pi$, let $x_1,\dots,x_n\in\ell_1$, and let $G$ be the resolvent matrix of $A$. Then
\[
\sum_{i,j}\pi_iG_{ij}\lnorm{x_i-x_j}^2\le28\sum_{i,j}\pi_i\Bigl(\frac1t\sum_{s=1}^t(A^s)_{ij}\Bigr)\lnorm{x_i-x_j}^2 .
\]
\end{lemma}
\begin{proof}
Let $D(s)=\sum_{i,j}\pi_i(A^s)_{ij}\lnorm{x_i-x_j}^2$.
\begin{itemize}
\item By \cref{lem:mc_resolvent_D_eq}, the left side equals $\sum_sp\,q^sD(s)$.
\item By \cref{lem:mc_cesaro_rhs_eq}, the sum on the right equals $\frac1t\sum_{s=1}^tD(s)$.
\item By \cref{lem:mc_D_props} and \cref{lem:mc_D_subadd}, $D$ satisfies the hypotheses of \cref{lem:mc_geom_cesaro}, which gives the bound. \qedhere
\end{itemize}
\end{proof}

\section{Proof of the theorem}

\begin{proof}[Proof of \cref{thm:main}]
Let $n,t,\pi,A,x_1,\dots,x_n$ be as in the theorem, and let $p=1/(t+1)$ and $q=t/(t+1)$.

\emph{Step 1: stationarity and cut representation.} By \cref{lem:mc_stationary}, $\pi A=\pi$. By \cref{lem:mc_cut_L1}, there are weights $w\in\R^{\Pn}$ with $w\ge0$ and a map $T:\R^{\Pn}\to\ell_1$ with $T(\chi_i)=x_i$, $\lnorm{x_i-x_j}=\wn{\chi_i-\chi_j}$, and $\lnorm{T(u)-T(u')}\le\wn{u-u'}$.

\emph{Step 2: the resolvent.} Let $G$ be the resolvent matrix, $G_{ij}=\sum_{s\ge0}p\,q^s(A^s)_{ij}$. By \cref{lem:mc_resolvent_basic}, $G_{ij}\ge0$, $\sum_jG_{ij}=1$ and $\sum_i\pi_iG_{ij}=\pi_j$. By \cref{lem:mc_resolvent_first_step}, $G=pI+qAG=pI+qGA$.

\emph{Step 3: the energy estimate on the cube.} Take $\iota=\Pn$ and $z_i=\chi_i\in\{0,1\}^{\Pn}$, and let $h_i=\sum_jG_{ij}\chi_j$ and $c_i=\Phi(h_i)$. By \cref{lem:mc_cube_energy},
\[
\sum_i\pi_i\wn{\chi_i-c_i}^2+t\sum_{i,j}\pi_ia_{ij}\wn{c_i-c_j}^2\le108\sum_{i,j}\pi_iG_{ij}\hn{\chi_i-\chi_j}^4 .
\]

\emph{Step 4: transfer to $\ell_1$.} The hypotheses of \cref{lem:mc_transfer} hold, since $\pi\ge0$ and $a_{ij}\ge0$. It yields points $y_i=T(\Phi(h_i))\in\ell_1$ with
\[
\sum_i\pi_i\lnorm{x_i-y_i}^2+t\sum_{i,j}\pi_ia_{ij}\lnorm{y_i-y_j}^2\le108\sum_{i,j}\pi_iG_{ij}\lnorm{x_i-x_j}^2 .
\]

\emph{Step 5: comparison.} By \cref{lem:mc_comparison}, which uses $t\ge1$ and $\pi A=\pi$, the right side is at most
\[
108\cdot28\sum_{i,j}\pi_i\Bigl(\frac1t\sum_{s=1}^t(A^s)_{ij}\Bigr)\lnorm{x_i-x_j}^2 .
\]
Since $108\cdot28=3024$, this proves the theorem.
\end{proof}

\paragraph{Relation to the source proof.}
The constructed points are the same as in the source: $y_i=T(\Phi(h_i))$, where $h_i$ is the resolvent average of the binary cut encodings. The constants are also the same: $108$ from the flatness lemma and $28$ from the geometric/Ces\`aro comparison. The formal argument deviates from the source in the following respects.

\begin{itemize}
\item \emph{No martingale.} The source realises the cost through a killed Markov chain and an $H$-valued martingale, and invokes conditional expectations and convergence theorems. The formal proof replaces this with deterministic algebra.
\begin{itemize}
\item The martingale property of one step is the first-step identity $h_j=p\,z_j+q\sum_ka_{jk}h_k$. This identity is what kills the linear term in \cref{lem:mc_one_step}.
\item The summation over time against the occupation measure of the killed chain becomes a weighting by $\pi_iG_{ij}$, which telescopes through $G=pI+qGA$ in \cref{lem:mc_aggregate}.
\item Discarding the initial potential corresponds to dropping the nonnegative terms $p\,\pi_iH_{ii}$.
\item The identity $B+tE_A$ becomes the final multiplication by $t+1$.
\end{itemize}
\item \emph{Comparison lemma.} The comparison of geometric and Ces\`aro averages is proved analytically for the sequence $D(s)$, and only for points of $\ell_1$. There is no random time $S$. The bounds $\mathbb E J^2\le3$ and $\mathbb P(R=r)\le2/t$ of the source appear as \cref{lem:mc_geom_sq_div} and \cref{lem:mc_geom_mod_coeff}.
\item \emph{Cut decomposition.} The decomposition uses all $2^n$ subsets of $[n]$ with weights $w_B\ge0$. Zero-weight cuts are not discarded, and there is no separate case for coinciding points. The cut vectors are obtained by sorting each coordinate and making a choice, rather than by the explicit min--max formula, though the resulting nested cut structure is the same.
\item \emph{Reversibility.} Reversibility is used only to derive stationarity (\cref{lem:mc_stationary}). Every later step uses only $\pi A=\pi$, in agreement with the remark in the source.
\end{itemize}

\appendix
\section{Correspondence with the formal proof}
Each lemma of the text and the Lean declaration that states it; \cref{thm:main} is \texttt{openaimath\_332\_metricmarkovcotypel1}.

\begin{longtable}{@{}lll@{}}
\toprule
Lemma & Lean declaration & Status \\
\midrule
\endhead
\cref{lem:mc_stationary} & \texttt{mc\_stationary} & verified \\
\cref{lem:mc_pow_stochastic} & \texttt{mc\_pow\_stochastic} & verified \\
\cref{lem:mc_pow_stationary} & \texttt{mc\_pow\_stationary} & verified \\
\cref{lem:mc_resolvent_summable} & \texttt{mc\_resolvent\_summable} & verified \\
\cref{lem:mc_resolvent_basic} & \texttt{mc\_resolvent\_basic} & verified \\
\cref{lem:mc_resolvent_first_step} & \texttt{mc\_resolvent\_first\_step} & verified \\
\cref{lem:mc_harmonic_point} & \texttt{mc\_harmonic\_point} & verified \\
\cref{lem:mc_phi_scalar} & \texttt{mc\_phi\_scalar} & verified \\
\cref{lem:mc_flat_increment} & \texttt{mc\_flat\_increment} & verified \\
\cref{lem:mc_flat} & \texttt{mc\_flat} & verified \\
\cref{lem:mc_one_step} & \texttt{mc\_one\_step} & verified \\
\cref{lem:mc_aggregate} & \texttt{mc\_aggregate} & verified \\
\cref{lem:mc_phi_vertex} & \texttt{mc\_phi\_vertex} & verified \\
\cref{lem:mc_cube_step} & \texttt{mc\_cube\_step} & verified \\
\cref{lem:mc_cube_energy_raw} & \texttt{mc\_cube\_energy\_raw} & verified \\
\cref{lem:mc_cube_energy} & \texttt{mc\_cube\_energy} & verified \\
\cref{lem:mc_cut_real_mono} & \texttt{mc\_cut\_real\_mono} & verified \\
\cref{lem:mc_cut_real_pointwise} & \texttt{mc\_cut\_real\_pointwise} & verified \\
\cref{lem:mc_cut_real} & \texttt{mc\_cut\_real} & verified \\
\cref{lem:mc_L1_norm} & \texttt{mc\_L1\_norm} & verified \\
\cref{lem:mc_L1_of_bound} & \texttt{mc\_L1\_of\_bound} & verified \\
\cref{lem:mc_L1_affine_apply} & \texttt{mc\_L1\_affine\_apply} & verified \\
\cref{lem:mc_L1_affine_contr} & \texttt{mc\_L1\_affine\_contr} & verified \\
\cref{lem:mc_L1_dist_eq} & \texttt{mc\_L1\_dist\_eq} & verified \\
\cref{lem:mc_cut_coords} & \texttt{mc\_cut\_coords} & verified \\
\cref{lem:mc_cut_L1} & \texttt{mc\_cut\_L1} & verified \\
\cref{lem:mc_transfer} & \texttt{mc\_transfer} & verified \\
\cref{lem:mc_minkowski} & \texttt{mc\_minkowski} & verified \\
\cref{lem:mc_weighted_tri} & \texttt{mc\_weighted\_tri} & verified \\
\cref{lem:mc_triple_mul} & \texttt{mc\_triple\_mul} & verified \\
\cref{lem:mc_triple_first} & \texttt{mc\_triple\_first} & verified \\
\cref{lem:mc_triple_second} & \texttt{mc\_triple\_second} & verified \\
\cref{lem:mc_D_subadd} & \texttt{mc\_D\_subadd} & verified \\
\cref{lem:mc_cesaro_Dt} & \texttt{mc\_cesaro\_Dt} & verified \\
\cref{lem:mc_cesaro_blocks} & \texttt{mc\_cesaro\_blocks} & verified \\
\cref{lem:mc_hasSum_sq_geom} & \texttt{mc\_hasSum\_sq\_geom} & verified \\
\cref{lem:mc_floor_div_sq_le} & \texttt{mc\_floor\_div\_sq\_le} & verified \\
\cref{lem:mc_geom_sq_div} & \texttt{mc\_geom\_sq\_div} & verified \\
\cref{lem:mc_geom_mod_coeff} & \texttt{mc\_geom\_mod\_coeff} & verified \\
\cref{lem:mc_geom_mod} & \texttt{mc\_geom\_mod} & verified \\
\cref{lem:mc_geom_summable_bdd} & \texttt{mc\_geom\_summable\_bdd} & verified \\
\cref{lem:mc_D_divmod} & \texttt{mc\_D\_divmod} & verified \\
\cref{lem:mc_geom_cesaro_pre} & \texttt{mc\_geom\_cesaro\_pre} & verified \\
\cref{lem:mc_range_le_Icc} & \texttt{mc\_range\_le\_Icc} & verified \\
\cref{lem:mc_geom_cesaro} & \texttt{mc\_geom\_cesaro} & verified \\
\cref{lem:mc_resolvent_D_eq} & \texttt{mc\_resolvent\_D\_eq} & verified \\
\cref{lem:mc_D_props} & \texttt{mc\_D\_props} & verified \\
\cref{lem:mc_cesaro_rhs_eq} & \texttt{mc\_cesaro\_rhs\_eq} & verified \\
\cref{lem:mc_comparison} & \texttt{mc\_comparison} & verified \\
\bottomrule
\end{longtable}
\end{document}
